%=======================================================================
%  CSE558 Data Science - Midterm Project Report  (course template, Monsoon 2026, IIIT-Delhi)
%  Project: Late deliveries and customer review scores on Olist
%  Compile with pdfLaTeX, twice, from the project root (figures are read from outputs/).
%=======================================================================
\documentclass[10pt,a4paper,twocolumn]{article}

%========================= EDIT THIS =========================
\newcommand{\projecttitle}{Impact of Delivery Delays on Customer Review Ratings:
  An Empirical Analysis of 95,296 Olist Marketplace Orders}
\newif\ifshowhints
\showhintsfalse      % \showhintsfalse hides every grey instruction line
% ---- Team details ----
\newcommand{\groupname}{Team tiriri}
\newcommand{\codelink}{\href{https://github.com/hansika24225/Data-Science-Midterm-project}{github.com/hansika24225/Data-Science-Midterm-project} (notebooks 01--04 and README)}
%=============================================================

%-------------------------------------------- fonts (no need to edit)
\usepackage{iftex}
\ifPDFTeX
  \usepackage[T1]{fontenc}
  \usepackage[utf8]{inputenc}
  \IfFileExists{carlito.sty}%
    {\usepackage[sfdefault]{carlito}}%              Calibri look
    {\usepackage[scaled=0.92]{helvet}\renewcommand{\familydefault}{\sfdefault}}%
  \usepackage{courier}
\else
  \usepackage[no-math]{fontspec}
  \setmainfont{Carlito}[Extension=.ttf, UprightFont=*-Regular, BoldFont=*-Bold,
    ItalicFont=*-Italic, BoldItalicFont=*-BoldItalic]
\fi

%----------------------------------------------------------- packages
\usepackage[margin=18mm,bottom=20mm,columnsep=7mm]{geometry}
\usepackage{amsmath}
\usepackage{xcolor}
\usepackage{graphicx}
\usepackage{tabularx}
\usepackage{array}
\usepackage{booktabs}
\usepackage{enumitem}
\usepackage{titlesec}
\usepackage{needspace}
\usepackage[font=small,labelfont=bf]{caption}
\usepackage[hidelinks]{hyperref}
\hypersetup{pdftitle={\projecttitle},
            pdfsubject={CSE558 Data Science, Monsoon 2026: Midterm Project Report}}

%------------------------------------------------------------- layout
\definecolor{eg}{HTML}{7396C4}         % worked example (light blue)
\definecolor{hintgrey}{gray}{0.58}     % instructions (light grey)

\titleformat{\section}{\normalfont\large\bfseries}{\thesection.}{0.5em}{}
\titlespacing*{\section}{0pt}{12pt}{3pt}
\titleformat{\subsection}{\normalfont\normalsize\bfseries}{\thesubsection}{0.5em}{}
\titlespacing*{\subsection}{0pt}{10pt}{2pt}
\newcommand{\sectionbreak}{\needspace{10\baselineskip}}
\newcommand{\subsectionbreak}{\needspace{10\baselineskip}}

\setlength{\parindent}{0pt}
\setlength{\parskip}{4pt}
\setlist{leftmargin=1.2em, itemsep=2pt, topsep=2pt, parsep=0pt}
\renewcommand{\arraystretch}{1.25}
\setlength{\tabcolsep}{4pt}
\newcolumntype{L}[1]{>{\raggedright\arraybackslash}p{#1}}
\newcolumntype{Y}{>{\raggedright\arraybackslash}X}
\newcolumntype{K}[1]{>{\raggedright\arraybackslash\bfseries\boldmath}p{#1}}  % bold labels

%--------------------------------------------------------- shorthands
\newcommand{\hint}[1]{\ifshowhints{\footnotesize\itshape\color{hintgrey}#1\par}\fi}
\newcommand{\eg}[1]{\textcolor{eg}{#1}}
\newcommand{\egblock}[1]{{\color{eg}#1\par}}
\newcommand{\field}[2]{\textbf{#1}~#2\par}
\newcommand{\tpart}[2]{\par\needspace{12\baselineskip}\smallskip\textbf{#1)~#2}\par}
\newcommand{\minihead}[1]{\par\needspace{8\baselineskip}\smallskip\textbf{#1}\par}
\newcommand{\code}[1]{\texttt{#1}}
\newcommand{\figfile}[2]{%
  \IfFileExists{#1}%
    {\includegraphics[width=\linewidth,height=#2,keepaspectratio]{#1}}%
    {\fbox{\parbox[c][#2][c]{\dimexpr\linewidth-2\fboxsep-2\fboxrule\relax}%
      {\centering\small\itshape\color{hintgrey}Figure: \detokenize{#1}}}}}
% Helper macros for p-values and figure directories
\newcommand{\pe}[1]{$10^{#1}$}
\newcommand{\Pthree}{outputs/phase3/figures}
\newcommand{\Pfour}{outputs/phase4/figures}
\renewcommand{\dbltopfraction}{0.85}\setcounter{dbltopnumber}{2}\renewcommand{\dblfloatpagefraction}{0.8}
\renewcommand{\topfraction}{0.85}\renewcommand{\textfraction}{0.1}

%=======================================================================
\begin{document}

\twocolumn[%
{\small\bfseries CSE558 Data Science\ \ \textbullet\ \ Midterm Project Report\ \ \textbullet\ \ Monsoon 2026\par}
\vspace{4pt}
{\LARGE\bfseries\raggedright\projecttitle\par}
\vspace{6pt}
\rule{\linewidth}{1pt}

{\small
\begin{tabularx}{\linewidth}{@{}K{2.4cm} Y@{}}
\toprule
Group name & \groupname \\
Dataset & kaggle.com/datasets/olistbr/brazilian-ecommerce (Brazilian E-Commerce Public Dataset by Olist) \\
Code & \codelink \\
\bottomrule
\end{tabularx}}

\vspace{2pt}
{\small
\begin{tabularx}{\linewidth}{@{}L{0.5cm} L{3.4cm} L{2.3cm} L{4.8cm} Y@{}}
\toprule
\textbf{\#} & \textbf{Name} & \textbf{Roll no.} & \textbf{Institute email} & \textbf{Contribution} \\
\midrule
1 & Divyanshi Gupta & SP26005 & divyanshi26005@iiitd.ac.in & Order-level aggregation, audit, 3M size check (1) \\
2 & Palak Khatri & SP26006 & palak26006@iiitd.ac.in & Cleaning decisions, before--after (2, 3) \\
3 & Hansika Singh & 2024225 & hansika24225@iiitd.ac.in & EDA, figures, problem statements (4) \\
4 & Rishit Sansanwal & 2024468 & rishit24468@iiitd.ac.in & Welch and z-tests, alternative tests, CIs (5) \\
5 & Aman Pratap Singh & SP26004 & aman26004@iiitd.ac.in & CMH and Spearman tests, sensitivity, limitations (5, 6) \\
\bottomrule
\end{tabularx}}
\vspace{10pt}
]

%-----------------------------------------------------------------------
\section{Dataset}
\field{Source:}{We use the ``Brazilian E-Commerce Public Dataset by Olist'' from Kaggle, which records real,
anonymised transactions on the Olist marketplace purchased September 2016 to August 2018. The dataset comprises
nine relational CSV files linked by order, customer, product and seller identifiers. Following order fulfillment,
buyers receive a satisfaction survey requesting a 1--5 star rating and optional written feedback. Our unit of
analysis is a single delivered order with exactly one review, giving a final primary sample of
\textbf{95,296 orders (6,353 late and 88,943 on time or early)}.}

\field{Raw structure and challenges:}{Order information is distributed across nine tables (orders 99,441\,$\times$\,8,
reviews 99,224\,$\times$\,7, items 112,650\,$\times$\,7, payments 103,886\,$\times$\,5, customers 99,441\,$\times$\,5,
products 32,951\,$\times$\,9, sellers 3,095\,$\times$\,4, translation 71\,$\times$\,2, and geolocation
1,000,163\,$\times$\,5, which is not required for order-level analysis). A single order can include multiple line
items, sellers and payment instalments. Furthermore, promised delivery dates are stored without time-of-day
timestamps, and customer comments consist of unstructured Portuguese text.}

{\small
\begin{tabularx}{\linewidth}{@{}L{3.0cm} Y Y@{}}
\toprule
 & \textbf{Joined table} & \textbf{After encoding} \\
\midrule
Instances (rows) & 95,296 & 95,296 \\
Attributes (columns) & 71 & 167 \\
Instances $\times$ attributes & 6,766,016 & 15,914,432 \\
\bottomrule
\end{tabularx}}

\field{Size check:}{$95{,}296\times167=15{,}914{,}432\ge3{,}000{,}000$, satisfying the course dimensionality
requirement. The 167 encoded columns combine 21 numeric or binary features with one-hot indicators for customer
state (27), purchase month (12) and purchase weekday (7), alongside multi-hot indicators for product category (74),
seller state (22) and payment method (4), while \code{order\_id} serves as the row index.}

\minihead{Key features}
{\footnotesize
\begin{tabularx}{\linewidth}{@{}L{3.5cm} L{1.25cm} Y@{}}
\toprule
\textbf{Feature} & \textbf{Type} & \textbf{Real-world meaning} \\
\midrule
\code{review\_score} & ordinal & Customer satisfaction rating (1--5 stars); primary response \\
\code{low\_score} & binary & Indicator equal to 1 if rating is 1 or 2 stars \\
\code{is\_late} & binary & Delivered on a later calendar day than promised; primary exposure \\
\code{delay\_calendar\_days} & discrete & Whole calendar days between promised and actual delivery date \\
\code{wait\_days} & continuous & Total waiting time from purchase to customer delivery (days) \\
\code{review\_before\_delivery} & binary & Indicator equal to 1 if review was answered before parcel arrival \\
\code{has\_comment} & binary & Indicator equal to 1 if review contains a non-blank text message \\
\code{customer\_state} & nominal & Buyer's Brazilian federative state (27 levels) \\
\code{freight\_ratio} & continuous & Total shipping cost divided by total item merchandise price \\
\bottomrule
\end{tabularx}}

%-----------------------------------------------------------------------
\begingroup\renewcommand{\sectionbreak}{}% the table below fills the rest of the column
\section{Data preparation}
\endgroup
{\footnotesize
\begin{tabularx}{\linewidth}{@{}L{2.0cm} L{2.4cm} Y@{}}
\toprule
\textbf{Issue found} & \textbf{Action taken} & \textbf{Rationale over alternatives} \\
\midrule
Orders not delivered or missing delivery date (2,971)
 & Excluded from sample
 & Actual waiting time and delivery lateness are undefined; imputing arrival dates would fabricate our primary exposure. \\
Orders lacking a review record (646)
 & Excluded from sample
 & The response variable is missing, leaving no outcome to evaluate. \\
Orders with multiple reviews (525 orders)
 & Excluded in primary sample
 & Multiple ratings create outcome ambiguity; retaining the latest review instead yields 95,824 orders with identical mean scores. \\
\bottomrule
\end{tabularx}}

{\footnotesize
\begin{tabularx}{\linewidth}{@{}L{2.0cm} L{2.4cm} Y@{}}
\toprule
\textbf{Issue found} & \textbf{Action taken} & \textbf{Rationale over alternatives} \\
\midrule
Multiple item and payment rows per order
 & Aggregated to order grain prior to joining; verified join cardinalities
 & Naive one-to-many joins duplicate multi-item orders and violate sample independence. \\
Promised date stored at midnight
 & Defined lateness by \emph{calendar day} rather than raw timestamp
 & Direct timestamp subtraction misclassifies 1,269 orders delivered later during the promised calendar day as late. \\
Survey creation vs.\ answer timestamp
 & Used answer timestamp (713 orders reclassified)
 & Only the submission timestamp reflects whether the buyer had received the parcel when rating. \\
Reviews answered prior to purchase (3)
 & Excluded from sample
 & Chronologically impossible timestamps where the order had not yet been placed. \\
\bottomrule
\end{tabularx}}

{\footnotesize
\begin{tabularx}{\linewidth}{@{}L{2.0cm} L{2.4cm} Y@{}}
\toprule
\textbf{Issue found} & \textbf{Action taken} & \textbf{Rationale over alternatives} \\
\midrule
Unprovided comment (59.4\%) and title (88.3\%)
 & Left as missing; created \code{has\_comment} and \code{has\_any\_text} indicators
 & Written feedback is optional in the survey, so imputing synthetic text is inappropriate. \\
Missing weight/size (16) and payment (1)
 & Preserved as \code{NaN} in main table; median-imputed only in encoded matrix
 & Our statistical tests do not depend on these columns, avoiding unnecessary row deletion. \\
Extreme durations: 290 waits $>$60 days, 341 delays $>$30 days
 & Retained with quality flags; parametric tests complemented by rank-based checks
 & These are genuine deliveries, not recording errors; waiting time ($\mu=12.51, \sigma=9.46$ days) is long-tailed yet well within the Markov ($P(X\ge30)=0.046\le0.417$) and Chebyshev ($P(|X-\mu|\ge2\sigma)=0.040\le0.250$) bounds. \\
Categorical attributes
 & One-hot (state, month, weekday) and multi-hot (category, seller state, payment)
 & Multi-item orders can span multiple categories, sellers and payment methods; $\sim$4,100 cities are too sparse to one-hot encode. \\
Duplicate and clustered records
 & No exact duplicates in the 8 active tables; flagged repeat buyers (5,163) and shared review IDs (867)
 & Repeat purchases represent real transactions whose clustering impact is evaluated via sensitivity subsets in Section 5. \\
\bottomrule
\end{tabularx}}

%-----------------------------------------------------------------------
\section{What changed after preparation}
{\footnotesize
\begin{tabularx}{\linewidth}{@{}L{2.9cm} L{1.6cm} Y@{}}
\toprule
\textbf{Measure} & \textbf{Before} & \textbf{After} \\
\midrule
Orders (rows) & 99,441 & 95,296 (one per order) \\
Late / on time or early & -- & 6,353 / 88,943 \\
Orders with several reviews & 525 & 0 \\
Impossible reviews (before purchase) & 3 & 0 \\
Orders misclassified by review timing & 713 & 0 \\
Columns of the analysis table & 55 & 71 (16 quality flags added) \\
Missing cells in the original 55 columns & 140,755 & 140,753 (kept missing) \\
Encoded matrix & -- & 95,296 $\times$ 167, no missing value \\
Original values changed & -- & 0 \\
\bottomrule
\end{tabularx}}

Our filtering funnel progressed from 99,441 raw orders $\to$ 96,478 delivered $\to$ 96,470 with a valid delivery
timestamp $\to$ 95,824 with at least one review $\to$ 95,299 with a single review $\to$ 95,296 after removing 3
pre-purchase review timestamps. Pre-aggregating line items and payments guaranteed unique \code{order\_id} records,
while using the survey answer timestamp corrected the timing classification of 713 orders—a critical step given that
70.0\% of late orders (compared to just 0.2\% of on-time orders) were reviewed prior to delivery.

%-----------------------------------------------------------------------
\section{Basic questions and problem statement}
\begin{figure*}[t]
  \centering
  \includegraphics[width=0.86\textwidth]{\Pthree/F01_score_distribution.png}
  \caption{\textbf{Late orders are heavily concentrated at 1-star ratings.} Distribution of review scores across all
  95,296 orders (Panel A) and stratified by delivery timeliness (Panel B: 6,353 late vs.\ 88,943 on time). Overall,
  59.2\% of orders receive 5 stars, whereas 62.4\% of late orders receive 1--2 stars compared with 9.2\% of on-time orders.}
\end{figure*}

\begin{figure*}[t]
  \centering
  \includegraphics[width=0.86\textwidth]{\Pthree/F09_review_timing_groups.png}
  \caption{\textbf{The rating penalty among late orders is driven primarily by reviews submitted before arrival.} Late
  orders reviewed before delivery ($n=4{,}450$) average 1.65 stars (80.7\% low scores); late orders reviewed after
  delivery ($n=1{,}903$) average 3.72 stars (19.4\% low), compared with 4.29 stars (9.2\% low) for on-time orders
  reviewed after delivery ($n=88{,}769$).}
\end{figure*}

\begin{figure*}[t]
  \centering
  \includegraphics[width=0.9\textwidth]{\Pthree/F08_low_score_by_wait_band.png}
  \caption{\textbf{Within identical waiting-time strata, late orders receive more low ratings, though the gap narrows
  substantially among post-delivery reviews.} Proportion of 1--2 star ratings across weekly waiting-time bands with
  95\% Wilson confidence intervals (strata with $<$30 orders omitted). Late and on-time distributions overlap primarily
  between 8 and 42 days.}
\end{figure*}

\minihead{Basic exploratory questions}
\begin{enumerate}[label=\arabic*.]
  \item \emph{How are customer ratings distributed overall?} Ratings are strongly left-skewed, with 59.2\% five-star
  and 9.7\% one-star reviews (Figure 1A).
  \item \emph{Do late deliveries receive lower ratings?} Yes: 62.4\% of late orders receive 1--2 stars compared with
  9.2\% of on-time orders (Figure 1B).
  \item \emph{When do buyers complete the review survey?} Overall, 70.0\% of late orders are reviewed before the
  package arrives (Figure 2).
  \item \emph{What do buyers discuss in written comments?} The phrase ``n\~ao recebi'' (``I did not receive'') appears
  in 24.4\% of late-order comments versus 1.9\% of on-time comments, and late orders include written messages
  58.9\% of the time compared to 39.3\% for on-time orders.
\end{enumerate}

\minihead{Novel problem statements}
\noindent\fbox{\begin{minipage}{\dimexpr\linewidth-2\fboxsep-2\fboxrule\relax}
\textbf{P1 (primary).} Conditioning on comparable actual waiting durations, are late deliveries still associated
with a higher incidence of 1--2 star ratings? (Evaluated in Test 3)

\smallskip
\textbf{P2.} Across the full marketplace sample, do late orders exhibit lower mean ratings, a higher proportion of
1--2 star reviews, and higher comment submission rates than on-time orders? (Evaluated in Tests 1, 2, and 5)

\smallskip
\textbf{P3.} Restricting to late orders, does greater delay severity correlate with lower review scores? (Evaluated in Test 4)
\end{minipage}}

\field{Why these problem statements are novel:}{Simple bivariate comparisons conflate \emph{missing the promised
delivery date} with \emph{longer total transit times}, since late orders take roughly three times longer to arrive
(median 31.1 vs.\ 9.8 days). Problem statement P1 disentangles expectation breach from elapsed duration by stratifying
across derived waiting-time bands, while our sensitivity framework re-evaluates every hypothesis on post-arrival
reviews to separate dissatisfaction with late arrival from frustration expressed while still waiting.}

%-----------------------------------------------------------------------
\section{Statistical inferences}
\field{Analytical samples:}{All 95,296 orders are evaluated in Tests 1, 2, and 5; Test 3 analyses the 68,080 orders
(4,988 late) falling within waiting-time bands containing at least 100 observations per group; Test 4 examines the
6,353 late orders.}

\field{Pre-specified testing protocol:}{We conduct five two-sided hypothesis tests while controlling the family-wise
error rate at $\alpha=0.05$ via the Holm--Bonferroni procedure. All group contrasts are defined as late minus on time.
Minimum practically meaningful effect thresholds were established following exploratory analysis (\code{analysis\_plan.json}),
and an observed effect is deemed practically significant when its entire 95\% confidence interval exceeds the threshold.}

%------------------------------------------------------------ Test 1
\subsection*{Test 1 (P2): Do late orders receive a lower mean review score?}
\tpart{a}{Hypotheses and pre-test specification}
{\small
\begin{tabularx}{\linewidth}{@{}K{3.4cm} Y@{}}
\toprule
$H_0$ & $\mu_{\text{late}}-\mu_{\text{on}}=0$ \\
$H_1$ & $\mu_{\text{late}}-\mu_{\text{on}}\neq0$ \\
Significance level & $\alpha=0.05$ (family-wise, Holm-adjusted) \\
Sidedness & Two-sided \\
Primary test & Two-sample Welch t-test \\
Smallest practical difference & 0.5 star \\
\bottomrule
\end{tabularx}}

Here $\mu$ denotes the population mean review score. A threshold of 0.5 stars represents half a rating unit on the
five-point scale, which is sufficient to visibly shift a product's displayed star average.

\tpart{b}{Alternative tests considered}
Comparing a numerical 1--5 response across two independent groups with unequal, unknown variances motivates Welch's
two-sample t-test. With sample sizes of 6,353 and 88,943, the sampling distribution of the mean is well-approximated
by a normal distribution via the Central Limit Theorem despite discrete skewness.
{\footnotesize
\begin{tabularx}{\linewidth}{@{}L{2.4cm} Y@{}}
\toprule
\textbf{Alternative considered} & \textbf{Reason for not selecting as primary} \\
\midrule
Two-sample z-test & Assumes known population standard deviations, whereas variances must be estimated from the sample. \\
Pooled (Student) t-test & Assumes homoscedasticity, which is rejected because the late group's variance is 1.88 times larger (Brown--Forsythe median Levene $F_{\text{BF}}=1{,}350.88, p=1.1\times10^{-293}$); running it yields $t=-131.58$ ($\text{df}=95{,}294$), same decision. \\
Mann--Whitney U test & Suitable for ordinal ratings and executed as a robustness check: $z=-96.08$, $p=\,$\pe{-2007}, where a random late order outscores a random on-time order 18.1\% of the time (95\% CI 17.5--18.7\%). Because it tests stochastic dominance rather than mean difference, it serves as a secondary check. \\
Paired t-test & Inapplicable because orders belong to two independent groups with no natural pairing. \\
\bottomrule
\end{tabularx}}

\tpart{c}{Results and interpretation}
{\small
\begin{tabularx}{\linewidth}{@{}K{3.4cm} Y@{}}
\toprule
Point estimate & $-2.02$ stars (2.27 late vs.\ 4.29 on time) \\
95\% confidence interval & $-2.06$ to $-1.98$ stars \\
Standardised effect size & Hedges' $g=-1.71$ \\
Test statistic (d.f.) & $t=-100.41$ (6,844) \\
p-value (raw / Holm) & \pe{-1347.7} / \pe{-1347.1} \\
Statistical decision & reject $H_0$ \\
\bottomrule
\end{tabularx}}

Late deliveries average more than two stars lower than on-time deliveries, and the entire 95\% confidence interval lies
roughly fourfold beyond the 0.5-star practical threshold. When restricted to reviews submitted after delivery, the mean
difference narrows to $-0.57$ ($-0.63$ to $-0.51$), which still exceeds the 0.5-star benchmark.

%------------------------------------------------------------ Test 2
\subsection*{Test 2 (P2): Do late orders have a higher proportion of 1--2 star ratings?}
\tpart{a}{Hypotheses and pre-test specification}
{\small
\begin{tabularx}{\linewidth}{@{}K{3.4cm} Y@{}}
\toprule
$H_0$ & $p_{\text{late}}-p_{\text{on}}=0$ \\
$H_1$ & $p_{\text{late}}-p_{\text{on}}\neq0$ \\
Significance level & $\alpha=0.05$ (family-wise, Holm-adjusted) \\
Sidedness & Two-sided \\
Primary test & Two-proportion z-test \\
Smallest practical difference & 5 percentage points \\
\bottomrule
\end{tabularx}}

Here $p$ denotes the proportion of orders receiving a 1--2 star rating. A 5 percentage point shift represents more than
one-third of the marketplace's baseline low-score rate of 12.8\%.

\tpart{b}{Alternative tests considered}
For a binary outcome across two independent samples where all expected contingency cell counts exceed 800, the normal
approximation of the two-proportion z-test is highly accurate and directly pairs with a signed Newcombe confidence interval.
{\footnotesize
\begin{tabularx}{\linewidth}{@{}L{2.4cm} Y@{}}
\toprule
\textbf{Alternative considered} & \textbf{Reason for not selecting as primary} \\
\midrule
Pearson $\chi^2$ test & Algebraically equivalent in two-sided form ($\chi^2=14{,}999.7=z^2$), but yields an unsigned statistic without a direct risk-difference interval. \\
$\chi^2$ with Yates' correction & Intended for small sample sizes; yields $\chi^2=14,995.0$, same decision. \\
Fisher's exact test & Computationally tailored for sparse tables; yields the same decision ($p<10^{-300}$, sample odds ratio 16.26). \\
\bottomrule
\end{tabularx}}

\tpart{c}{Results and interpretation}
{\small
\begin{tabularx}{\linewidth}{@{}K{3.4cm} Y@{}}
\toprule
Point estimate & $+53.1$ pp (62.4\% late vs.\ 9.2\% on time) \\
95\% confidence interval & 51.9 to 54.3 pp \\
Relative effect size & Risk ratio 6.74 (6.56 to 6.94) \\
Test statistic & $z=122.47$ \\
p-value (raw / Holm) & \pe{-3259.3} / \pe{-3258.6} \\
Statistical decision & reject $H_0$ \\
\bottomrule
\end{tabularx}}

Approximately 62 out of every 100 late orders receive 1--2 stars compared with only 9 out of 100 on-time orders, far
surpassing the 5 percentage point practical threshold. Among post-delivery reviews, the difference drops to $+10.2$ pp
(8.4 to 12.0; risk ratio 2.10), remaining practically meaningful.

%------------------------------------------------------------ Test 3
\subsection*{Test 3 (P1, primary): Does the low-score association persist across comparable waiting times?}
\tpart{a}{Hypotheses and pre-test specification}
{\small
\begin{tabularx}{\linewidth}{@{}K{3.4cm} Y@{}}
\toprule
$H_0$ & Common odds ratio $=1$ across waiting-time strata \\
$H_1$ & Common odds ratio $\neq1$ \\
Significance level & $\alpha=0.05$ (family-wise, Holm-adjusted) \\
Sidedness & Two-sided \\
Primary test & Cochran--Mantel--Haenszel (CMH) test \\
Smallest practical difference & Common odds ratio 1.5 \\
\bottomrule
\end{tabularx}}

The stratum-specific odds ratio compares the odds of a 1--2 star rating between late and on-time deliveries \emph{within
the same weekly waiting-time band}. Following our pre-specified eligibility rule requiring at least 100 late and 100
on-time orders per stratum, the extreme bands of 0--7 days (58 late), 43--60 days (79 on time), and 60+ days (3 on time)
are excluded, leaving 68,080 orders across the 8--42 day overlap region (with 27,216 non-overlapping orders excluded).

\tpart{b}{Alternative tests considered}
The CMH procedure tests conditional independence across $2\times2\times K$ strata non-parametrically and provides a
pooled Mantel--Haenszel odds ratio; all expected cell counts within the four eligible strata are at least 28.
{\footnotesize
\begin{tabularx}{\linewidth}{@{}L{2.4cm} Y@{}}
\toprule
\textbf{Alternative considered} & \textbf{Reason for not selecting as primary} \\
\midrule
Unstratified two-proportion z-test & Addresses P2 rather than P1 by confounding expectation breach with total waiting duration (crude odds ratio 12.86). \\
Logistic regression with stratum indicators & Imposes parametric logit linearity and relies on Wald asymptotics; yields a similar adjusted odds ratio of 6.81 (6.29 to 7.38), same decision. \\
Continuity-corrected CMH & Continuity adjustment is only needed for sparse strata; yields $\chi^2=2468.77$, same decision. \\
\bottomrule
\end{tabularx}}

\tpart{c}{Results and interpretation}
{\small
\begin{tabularx}{\linewidth}{@{}K{3.4cm} Y@{}}
\toprule
Point estimate & Common odds ratio 7.03 \\
95\% confidence interval & 6.48 to 7.63 \\
Test statistic (d.f.) & $\chi^2=2470.98$ (1) \\
p-value (raw / Holm) & \pe{-538.4} / \pe{-537.9} \\
Statistical decision & reject $H_0$ \\
Post-delivery subset (S1) & OR 1.41 (1.24 to 1.60), $p=2\times10^{-7}$ \\
\bottomrule
\end{tabularx}}

{\footnotesize
\begin{tabularx}{\linewidth}{@{}L{1.3cm} L{1.7cm} Y L{1.9cm}@{}}
\toprule
\textbf{Band (days)} & \textbf{Late / on time} & \textbf{1--2 stars, late vs.\ on time} & \textbf{Odds ratio (95\% CI)} \\
\midrule
8--14 & 313 / 39,468 & 32.6\% vs.\ 8.8\% & 5.04 (3.97--6.40) \\
15--21 & 526 / 16,927 & 43.2\% vs.\ 10.7\% & 6.36 (5.32--7.61) \\
22--28 & 1,490 / 5,295 & 50.3\% vs.\ 14.2\% & 6.10 (5.37--6.93) \\
29--42 & 2,659 / 1,402 & 69.6\% vs.\ 19.4\% & 9.50 (8.13--11.10) \\
\bottomrule
\end{tabularx}}

Conditioning on waiting duration reduces the crude odds ratio (12.86) by nearly half, suggesting that elapsed wait time
accounts for part of the raw association, yet late orders still face a sevenfold increase in the odds of a low rating—well above our
1.5 threshold. Although Breslow--Day--Tarone heterogeneity is statistically detectable ($p\approx5\times10^{-6}$), the
ratio of maximum to minimum stratum ORs remains below 2, supporting a pooled summary. To address residual within-band
median wait differences of 1.0--3.1 days, repeating CMH with finer half-week strata yields a consistent OR of 6.55
(6.02 to 7.12). Crucially, restricting to post-delivery reviews reduces the pooled odds ratio to \textbf{1.41}, falling
below the 1.5 practical threshold and showing that the bulk of the primary effect stems from buyers rating while
still awaiting overdue packages.

\begin{figure*}[t]
  \centering
  \includegraphics[width=0.78\textwidth]{\Pfour/P4_F2_stratified_odds_ratios.png}
  \caption{\textbf{Test 3: Late deliveries exhibit substantially higher low-score odds within every waiting-time band,
  whereas the post-delivery effect is modest.} Stratum-specific odds ratios with 95\% CIs (log scale) and pooled
  Mantel--Haenszel estimates (diamonds) for all eligible orders (black: 68,080 total; 4,988 late) versus reviews
  submitted after delivery (violet: 64,737 total; 1,811 late), relative to the 1.5 practical threshold.}
\end{figure*}

%------------------------------------------------------------ Test 4
\subsection*{Test 4 (P3): Among late orders, does longer delay duration correlate with lower ratings?}
\tpart{a}{Hypotheses and pre-test specification}
{\small
\begin{tabularx}{\linewidth}{@{}K{3.4cm} Y@{}}
\toprule
$H_0$ & $\rho_S=0$ between calendar days late and review score \\
$H_1$ & $\rho_S\neq0$ \\
Significance level & $\alpha=0.05$ (family-wise, Holm-adjusted) \\
Sidedness & Two-sided \\
Primary test & Spearman rank correlation \\
Smallest practical difference & $|\rho_S|=0.10$ \\
\bottomrule
\end{tabularx}}

Delay severity is measured via \code{delay\_calendar\_days} ($\ge1$ across all late deliveries), and $|\rho_S|=0.10$
corresponds to the standard benchmark for a non-trivial rank correlation.

\tpart{b}{Alternative tests considered}
Because delay durations are heavily right-skewed and star ratings are ordinal, Spearman's rank correlation captures
monotonic association without requiring bivariate normality; its 95\% CI is computed via BCa bootstrap (5,000 resamples,
seed 558).
{\footnotesize
\begin{tabularx}{\linewidth}{@{}L{2.4cm} Y@{}}
\toprule
\textbf{Alternative considered} & \textbf{Reason for not selecting as primary} \\
\midrule
Pearson linear correlation & Assumes interval linearity and is sensitive to extreme delays reaching 188 days, attenuating the estimate to $r=-0.15$ ($-0.17$ to $-0.12$). \\
Kendall's $\tau_b$ & A valid concordance alternative that yields $\tau_b=-0.31$ ($p=5.8\times10^{-226}$, same decision), but operates on a compressed numerical scale relative to $\rho_S$. \\
\bottomrule
\end{tabularx}}

\tpart{c}{Results and interpretation}
{\small
\begin{tabularx}{\linewidth}{@{}K{3.4cm} Y@{}}
\toprule
Point estimate & $\rho_S=-0.40$ ($n=6{,}353$ late orders) \\
95\% confidence interval & $-0.42$ to $-0.38$ \\
Test statistic (d.f.) & $t=-35.04$ (6,351) \\
p-value (raw / Holm) & $5.1\times10^{-246}$ / $1.0\times10^{-245}$ \\
Statistical decision & reject $H_0$ \\
Post-delivery subset (S1) & $-0.03$ ($-0.08$ to $0.01$), $p=0.13$ \\
\bottomrule
\end{tabularx}}

Across all late orders, delay length exhibits a moderate negative rank correlation with rating that clearly exceeds the
0.10 benchmark. However, Figure 5 demonstrates that this downward slope is driven by pre-arrival reviews: among
the $n=1{,}903$ late orders reviewed after delivery, delay duration shows no meaningful association with rating ($\rho_S=-0.03$),
with its entire 95\% CI ($[-0.08, 0.01]$) contained strictly inside the $[-0.10, +0.10]$ practical equivalence zone.

\begin{figure*}[t]
  \centering
  \includegraphics[width=0.74\textwidth]{\Pfour/P4_F3_delay_vs_score.png}
  \caption{\textbf{Test 4: Review scores decline with additional days of delay only among surveys answered prior to
  delivery.} Mean review score across calendar-day delay bins with 95\% CIs (bins with $<$30 orders omitted), contrasting
  all late orders ($\rho_S=-0.40, n=6{,}353$) with late orders reviewed after arrival ($\rho_S=-0.03, n=1{,}903$).}
\end{figure*}

%------------------------------------------------------------ Test 5
\subsection*{Test 5 (P2): Are late orders more likely to include a written comment?}
\tpart{a}{Hypotheses and pre-test specification}
{\small
\begin{tabularx}{\linewidth}{@{}K{3.4cm} Y@{}}
\toprule
$H_0$ & $c_{\text{late}}-c_{\text{on}}=0$ \\
$H_1$ & $c_{\text{late}}-c_{\text{on}}\neq0$ \\
Significance level & $\alpha=0.05$ (family-wise, Holm-adjusted) \\
Sidedness & Two-sided \\
Primary test & Two-proportion z-test \\
Smallest practical difference & 5 percentage points \\
\bottomrule
\end{tabularx}}

Here $c$ denotes the proportion of orders containing a non-blank comment message (\code{has\_comment}), capturing active
customer feedback engagement.

\tpart{b}{Alternative tests considered}
As in Test 2, we compare a binary indicator across two large independent samples where the minimum expected cell count is 2,580.
{\footnotesize
\begin{tabularx}{\linewidth}{@{}L{2.4cm} Y@{}}
\toprule
\textbf{Alternative considered} & \textbf{Reason for not selecting as primary} \\
\midrule
$\chi^2$ with Yates' correction & Unnecessary for large expected counts; yields $\chi^2=947.1$, $p=5.6\times10^{-208}$, same decision. \\
Fisher's exact test & Designed for small samples; yields odds ratio 2.22, $p=1.4\times10^{-203}$, same decision. \\
\code{has\_any\_text} (message or title) & Evaluates a broader outcome definition; as a robustness check it yields $+19.0$ pp (17.7 to 20.2). \\
\bottomrule
\end{tabularx}}

\tpart{c}{Results and interpretation}
{\small
\begin{tabularx}{\linewidth}{@{}K{3.4cm} Y@{}}
\toprule
Point estimate & $+19.6$ pp (58.9\% late vs.\ 39.3\% on time) \\
95\% confidence interval & 18.4 to 20.9 pp \\
Test statistic & $z=30.79$ \\
p-value (raw / Holm) & $3.7\times10^{-208}$ / $3.7\times10^{-208}$ \\
Statistical decision & reject $H_0$ \\
Post-delivery subset (S1) & $+3.8$ pp (1.6 to 6.1) \\
\bottomrule
\end{tabularx}}

While the overall $+19.6$ pp increase in comment propensity is substantial, restricting to post-delivery reviews reduces
the difference to $+3.8$ pp (below our 5 pp threshold), indicating that heightened comment activity mainly comes from
buyers who had not yet received their packages.

%------------------------------------------------------------ Holm and sensitivity
\subsection*{Multiple testing control and sensitivity analysis}
\textbf{Holm--Bonferroni procedure.} Running five tests at $\alpha=0.05$ without adjustment can raise the family-wise
Type I error rate to as much as 22.6\% (if the tests were independent). Applying Holm's step-down adjustment (multiplying ordered p-values by $5, 4, 3, 2, 1$)
strictly bounds the family-wise error rate at 5\%, with all five null hypotheses rejected. Where test statistics produce
p-values below IEEE-754 double precision ($10^{-300}$), exact magnitudes are computed in $\log_{10}p$ space rather than
reported as zero.

\textbf{Sensitivity subsets.} Pre-specified robustness checks:
{\footnotesize
\begin{tabularx}{\linewidth}{@{}L{1.75cm} Y Y Y Y Y@{}}
\toprule
\textbf{Sample} & \textbf{T1 stars} & \textbf{T2 pp} & \textbf{T3 OR} & \textbf{T4 $\rho_S$} & \textbf{T5 pp} \\
\midrule
Primary & $-$2.02 & 53.1 & 7.03 & $-$0.40 & 19.6 \\
S1 after delivery & $-$0.57 & 10.2 & 1.41 & $-$0.03 & 3.8 \\
S2 one per customer & $-$2.02 & 53.3 & 7.04 & $-$0.40 & 19.7 \\
S3 no shared reviews & $-$2.02 & 53.2 & 7.11 & $-$0.41 & 19.6 \\
S4 all three & $-$0.57 & 10.0 & 1.38 & $-$0.04 & 3.5 \\
\bottomrule
\end{tabularx}}

Subsample S1 retains the 90,672 orders reviewed after delivery (1,903 late); S2 retains only the first order per unique
buyer across 92,583 customers; S3 removes 867 orders sharing a review ID; and S4 combines all three restrictions
(87,736 orders; 1,845 late). Eliminating repeat buyers (S2) or shared review IDs (S3) leaves all estimates virtually
unchanged, whereas conditioning on post-delivery reviews (S1, S4) substantially attenuates every effect size (see
forest plot P4\_F1 in \code{outputs/phase4/figures}). Using continuous timestamp delays in Test 4 similarly yields $\rho_S=-0.405$.

\minihead{Summary conclusion}
\noindent\fbox{\begin{minipage}{\dimexpr\linewidth-2\fboxsep-2\fboxrule\relax}
Late delivery is \textbf{strongly associated} with lower customer ratings ($-2.02$ stars and a 6.74-fold higher risk of
1--2 star reviews), and this relationship \textbf{persists across comparable waiting times} (common OR = 7.03 across
8--42 days). However, the penalty is \textbf{substantially smaller among reviews submitted after parcel arrival} ($-0.57$
stars; common OR = 1.41, below our 1.5 practical threshold), where \textbf{additional days of delay no longer correlate}
with rating ($\rho_S=-0.03$). Thus, customer dissatisfaction is heavily concentrated among buyers prompted to rate while
still awaiting overdue orders.
\end{minipage}}

%-----------------------------------------------------------------------
\section{Limitations}
\begin{itemize}
  \item \textbf{Observational association vs.\ causality.} Delivery timeliness is not randomly assigned. Late and
  on-time orders differ systematically across buyer states (where state-level late rates range from 2.8\% to 20.8\%),
  merchant profiles, product categories, and seasonal volume spikes; Test 3 controls only for total waiting time.
  \item \textbf{Survey timing selection.} Overall, 70.0\% of late orders were rated prior to physical delivery. These
  responses capture real frustration during an overdue wait rather than post-receipt product evaluation, and buyers who
  wait until delivery to respond may differ systematically.
  \item \textbf{Partial waiting-time overlap.} Late and on-time distributions overlap sufficiently only within the
  8--42 day window, excluding 27,216 extreme orders from Test 3, and mild within-stratum wait differences remain.
  \item \textbf{Non-independence of observations.} Repeat purchases (5,163 orders), shared review IDs (867 orders),
  and heavy seller concentration (Gini $=0.782$, top 5\% of 2,963 sellers fulfilling 52.9\% of volume) introduce
  clustering. Subsets S2 and S3 show that repeat buyers and shared reviews do not change our conclusions; seller-level
  clustering is not adjusted for, so p-values may be somewhat optimistic.
  \item \textbf{Scale ordinality.} Parametric mean differences assume equal intervals between star ratings, which we
  mitigate by complementing t-tests with binary low-score proportions and non-parametric rank tests.
  \item \textbf{Sample scope.} Findings apply to completed, single-review deliveries on Olist between 2016 and 2018,
  excluding 2,971 undelivered orders and 525 multi-review orders.
\end{itemize}

\end{document}
